The comparative analysis of the simulation trajectories unequivocally demonstrates the superior efficacy of the Adaptive Heuristic-Guided Deep Reinforcement Learning (HuDRL) Strategy over all static, equal-distribution policy interventions. As detailed in Table \ref{tab:policy_comparison}, traditional baseline and single-instrument strategies systematically fail to generate or sustain municipality-wide behavioral change \citep{Villanueva2021}. This failure is not fundamentally due to a lack of total municipal funding, but rather the mathematical and behavioral realities of resource dilution when applied across heterogeneous barangays \citep{Alphonso2024, Badua2022}.

\begin{table}[ht]
    \centering
    \caption{Comparative Global Compliance Rates across Policy Regimes (Years 1--3)}
    \label{tab:policy_comparison}
    \small
    \renewcommand{\arraystretch}{1.2}
    \begin{tabular}{l c c c c}
        \toprule
        \textit{Policy Regime} & \textit{Q1 (Year 1)} & \textit{Q4 (Year 1)} & \textit{Q8 (Year 2)} & \textit{Q12 (Year 3)} \\
        \midrule
        Status Quo (Baseline) & 12.50\% & 12.72\% & 11.50\% & 10.63\% \\
        Pure Enforcement & 12.50\% & 29.50\% & 29.23\% & 29.60\% \\
        Pure Incentives & 12.50\% & 28.62\% & 28.64\% & 32.20\% \\
        \textbf{Adaptive (HuDRL)} & \textbf{12.50\%} & \textbf{73.65\%} & \textbf{97.84\%} & \textbf{98.23\%} \\
        \bottomrule
    \end{tabular}
    \vspace{0.2cm}
    \begin{minipage}{0.9\textwidth}
        \footnotesize
        \textit{Note: Global Compliance is calculated as the true population-weighted mean of the segregation rates across the seven constituent barangays, accurately reflecting the impact of high-density areas. While static policies stagnate due to budget dilution, the Adaptive Strategy immediately outpaces baselines by Quarter 4 through aggressive, targeted resource allocation.}
    \end{minipage}
\end{table}

\subsection{The Dilution Effect and the Failure of the Status Quo}
The Status Quo approach, heavily reliant on Information, Education, and Communication (IEC) campaigns distributed equally across the municipality, experienced a steady decay, ending at a dismal terminal compliance of 10.63\%. This trajectory computationally validates the limitations of the ``Value-Action Gap''---also referred to as the intention-action gap in environmental psychology \citep{Geiger2021, Liao2024}. While IEC campaigns may marginally increase the environmental attitude ($w_a$) of the household agents, passive education fails to build habit formation when the physical friction of segregation remains unaddressed \citep{Moeini2023, Yazawa2025}. By spreading limited funds equally, the Local Government Unit (LGU) dilutes its capacity, ensuring that the critical density of policing or the minimum per-capita incentive threshold is never met in any single barangay \citep{Alphonso2024}. 

\subsection{The Plateau of Single-Instrument Regimes}
Similarly, the simulation revealed the inherent behavioral ceilings of single-instrument regimes. Both the Pure Enforcement (29.60\%) and Pure Incentives (32.20\%) strategies plateaued early in the first year and stagnated through Year 3. 

In the Pure Enforcement scenario, the uniform distribution of penal resources meant that dense urban centers like Barangay Poblacion simply overwhelmed the available enforcers, a phenomenon consistent with the \textit{Urban Resource Trap} \citep{Villanueva2021}. The probability of a non-compliant household being caught remained statistically insignificant, neutralizing the deterrent effect of the policy \citep{Badua2022}. Conversely, the Pure Incentives scenario failed due to the diminishing per-capita value of rewards \citep{MuZhang2021}. When a fixed municipal budget is divided equally among a large population, the individual payout drops below the household's ``Cost of Effort'' ($c_{\text{effort}}$) threshold, rendering the financial motivation insufficient to overcome the logistical inconvenience of waste sorting \citep{Chen2023, Corcoran2020}.

\subsection{Sequential Saturation: The HuDRL Advantage}
In stark contrast, the HuDRL Agent successfully navigated these socio-economic trade-offs by dynamically shifting resources, achieving an exceptional terminal Global Compliance Rate of 98.23\% by Quarter 12 \citep{Dey2025, ZhengAI2022}. This represents a massive deviation from the stagnant, linear trajectories of the static models. 

Rather than suffering from a slow exploration phase, the adaptive strategy broke the municipality's inertia almost immediately, surging to 73.65\% compliance by the end of Year 1. This overwhelming success is attributed directly to the agent's discovery of the ``Sequential Saturation'' logic. The AI recognized that overcoming the urban density barrier required abandoning the heuristic of simultaneous, equal distribution. Instead, the agent utilized sequential dominance---concentrating its budget to artificially induce a behavioral tipping point in specific barangays \citep{Centola2018, Nyborg2024}, locking in compliance through social norms \citep{Andre2021}, and subsequently reallocating those funds to besiege the next resistant population center \citep{Nishimura2022}. This dynamic phasing proved that in resource-constrained LGUs, the timing and concentration of capital are far more critical than the absolute budget size \citep{Jimenez2025}.